\section{History of the Topic}

\subsection{The Government Monopoly Era (1957–1980s)}

The exploration of outer space began as an exclusively state-driven enterprise during the Cold War, with the United States and the Soviet Union competing for dominance beyond Earth’s atmosphere. In 1967, the international community adopted the Treaty on Principles Governing the Activities of States in the Exploration and Use of Outer Space (the Outer Space Treaty), which established outer space as “the province of all mankind,” free for exploration and use by all nations, and explicitly banned national appropriation of the Moon or other celestial bodies by sovereignty claims.\footnote{United Nations Office for Outer Space Affairs. (n.d.). \emph{Treaty on Principles Governing the Activities of States in the Exploration and Use of Outer Space, Including the Moon and Other Celestial Bodies}. United Nations. \url{https://www.unoosa.org/oosa/en/ourwork/spacelaw/treaties/introouterspacetreaty.html}} During this period, NASA maintained a monopoly on launching spacecraft and sending humans into space, though it always worked with private aerospace contractors—such as General Dynamics, Boeing, and Lockheed Martin—on design, innovation, and feasibility studies.\footnote{Dagnino, G., \& Orrù, E. (2023). Outsourcing the American space dream: SpaceX and the narrative of privatized space exploration. \emph{Geopolitics}, ahead-of-print. \url{https://www.diva-portal.org/smash/get/diva2:2048145/FULLTEXT01.pdf}} The legal and institutional framework created during this era set the boundaries within which all future privatization efforts would have to operate.

\subsection{Early Privatization (1984–2000s)}

The shift toward commercial involvement in space began in the early 1980s when Space Services Inc. developed the world’s first privately funded rocket, the Conestoga 1, which successfully flew in September 1982.\footnote{Harrison, T. (2023). The commercial space imperative. Center for Strategic and International Studies. \url{https://www.csis.org/analysis/commercial-space-imperative}} The approval procedures for that launch led to legislation making it easier for companies to pursue commercial launch activities.\footnote{Federal Aviation Administration. (n.d.). \emph{Origins of the commercial space industry}. U.S. Department of Transportation. \url{https://www.faa.gov/sites/faa.gov/files/about/history/milestones/Commercial_Space_Industry.pdf}} On January 25, 1984, President Reagan stated in his State of the Union address that the market for space transportation could surpass the nation’s capacity to develop it, and that private companies must have access to launch services.\footnote{U.S. Senate Committee on Commerce, Science, and Transportation. (2015). \emph{U.S. Commercial Space Launch Competitiveness Act} (Senate Report 114-88). U.S. Government Publishing Office. \url{https://www.govinfo.gov/content/pkg/CRPT-114srpt88/html/CRPT-114srpt88.htm}} Later that year, Congress passed the Commercial Space Launch Act of 1984 (P.L. 98-575), which recognized the U.S. private sector’s capability to develop commercial launch vehicles and assigned the Department of Transportation the duties of overseeing commercial launches and issuing licenses.\footnote{Commercial Space Launch Act of 1984, Pub. L. No. 98-575, 98 Stat. 3055. \url{https://www.govtrack.us/congress/bills/98/hr3942/text}} Reagan’s policy succeeded in turning space launch over to the commercial world: by 1998, the number of commercial launches outnumbered government launches for the first time.

\subsection{NASA’s Commercial Transition (2000s–2010s)}

A paradigm shift occurred in 2006 when NASA initiated the Commercial Orbital Transportation Services (COTS) program, selecting private space companies to develop and demonstrate new capabilities for resupplying the International Space Station (ISS), with funding tied to milestone achievement rather than traditional cost-plus contracts.\footnote{Handberg, R. (2015). Launching commercial space: NASA, cargo, and policy innovation. \emph{Space Policy}, \emph{34}, 23–31. \url{https://doi.org/10.1016/j.spacepol.2015.08.001}} NASA originally awarded COTS Phase 1 funding to SpaceX and Rocketplane Kistler in 2006; however, Rocketplane Kistler failed to meet private fundraising milestones, and NASA terminated their contract in 2007.\footnote{The Breakthrough Institute. (n.d.). Commercial spaceflight. \url{https://thebreakthrough.org/issues/energy/commercial-spaceflight}} SpaceX went from announcing a new rocket to successfully joining with the ISS in 2012, reducing the initial cost of payload delivery to roughly \$9,000 per pound, comparable to what NASA’s Space Shuttle had achieved. In 2010, NASA invested nearly \$50 million in the Commercial Crew Development program (CCDev1) to stimulate private sector efforts in crew transportation, and in 2014 selected Boeing and SpaceX under the Commercial Crew Transportation Capability (CCtCap) contracts to transport crews to the ISS.\footnote{National Aeronautics and Space Administration. (n.d.). \emph{Commercial Crew Program essentials}. NASA. \url{https://www.nasa.gov/humans-in-space/commercial-space/commercial-crew-program/commercial-crew-program-essentials/}}

\subsection{The Current Commercial Boom and the Colonization Question}

Today, the commercial space sector has expanded dramatically, with billions in private investment across more than a hundred companies, massive reductions in government launch costs, and the emergence of space tourism as a viable market.\footnote{Harrison, T. (2023). The commercial space imperative. Center for Strategic and International Studies. \url{https://www.csis.org/analysis/commercial-space-imperative}} Fully private crews now fly on commercially owned and operated human space systems for both suborbital and orbital flights.\footnote{National Aeronautics and Space Administration. (n.d.). How NASA's work led to commercial spaceflight revolution. NASA. \url{https://www.nasa.gov/humans-in-space/commercial-space/how-nasas-work-led-to-commercial-spaceflight-revolution/}} However, the question of space colonization raises profound legal and ethical challenges. The Artemis Accords, presented in October 2020 by a coalition of nations, aim to operationalize the norms of the Outer Space Treaty for lunar exploration and resource extraction, asserting that the extraction of space resources does not inherently constitute national appropriation under Article II of the Treaty.\footnote{National Aeronautics and Space Administration. (2020). \emph{The Artemis Accords: Principles for cooperation in the civil exploration and use of the Moon, Mars, comets, and asteroids for peaceful purposes}. \url{https://www.nasa.gov/artemis-accords/}} Critics warn that this framework represents an attempt to legitimize deviation from the Outer Space Treaty and may not lead to a fair distribution of resources.\footnote{Deplano, R. (2021). The Artemis Accords: Evolution or revolution in international space law? \emph{International and Comparative Law Quarterly}, \emph{72}(3), 671–714. \url{https://doi.org/10.1017/S0020589323000234}} For the House of Representatives, the key legislative questions include how to balance encouraging private investment with ensuring compliance with international obligations, protecting public safety, and addressing whether commercial colonization efforts align with the principle that outer space remains the province of all mankind.
